% MG06. En exc:paley-codigo, tras definir C y chi, antes de la frase
% «Las identidades de la suma cuadrática...» y de reducir módulo tres.
La identidad cuadrática de esta coordenatización se obtiene mediante una
correlación finita. Se conserva el orden
\((\infty,0,1,2,3,4)\), junto con la identificación de las seis
posiciones y el carácter \(\chi\) ya declarados.

\begin{lemma}[Correlación del carácter cuadrático en cinco posiciones]
\label{mg06:correlacion}
Para \(a,b\in\mathbb F_5\) con \(a\ne b\), se cumple
\begin{equation}
 \sum_{t\in\mathbb F_5}\chi(t-a)\chi(t-b)=-1.
 \label{mg06:identidad-correlacion}
\end{equation}
\end{lemma}
\begin{proof}
La sustitución \(u=(t-a)/(b-a)\) es biyectiva, pues \(b-a\) es
invertible en \(\mathbb F_5\). Además,
\[
 (t-a)(t-b)=(b-a)^2u(u-1),\qquad
 \chi((b-a)^2)=1.
\]
La multiplicatividad de \(\chi\), incluida su extensión por cero,
reduce la suma a \(\sum_u\chi(u(u-1))\). Las cinco evaluaciones son
\begin{equation}
\begin{array}{c|rrrrr}
 u&0&1&2&3&4\\\hline
 u(u-1)\pmod5&0&0&2&1&2\\
 \chi(u(u-1))&0&0&-1&1&-1
\end{array}
\label{mg06:cinco-evaluaciones}
\end{equation}
Su suma es \(-1\), lo que demuestra
\eqref{mg06:identidad-correlacion}.
\end{proof}

\begin{proposition}[Ortogonalidad de la coordenatización de seis posiciones]
\label{mg06:conferencia}
La matriz \(C\) satisface
\begin{equation}
 C=C^{\mathsf T},\qquad CC^{\mathsf T}=5I_6,
 \qquad C^2=5I_6.
 \label{mg06:identidad-conferencia}
\end{equation}
\end{proposition}
\begin{proof}
Como \(\chi(-1)=\chi(4)=1\), se tiene
\(\chi(b-a)=\chi(a-b)\); las entradas que involucran
\(\infty\) también son simétricas. Cada fila contiene una entrada
nula y cinco entradas de módulo uno, de modo que su norma cuadrada
es cinco. Para dos filas finitas distintas, el producto escalar
recibe \(+1\) de la coordenada \(\infty\) y \(-1\) de las cinco
coordenadas finitas por \eqref{mg06:identidad-correlacion}; su suma
es cero. El producto escalar de la fila infinita con una fila finita
es \(\sum_{t\in\mathbb F_5}\chi(t)=0+1-1-1+1=0\).
Quedan evaluadas todas las entradas de \(CC^{\mathsf T}\).
La simetría convierte esta identidad en \(C^2=5I_6\).
\end{proof}

El argumento conserva el orden de la coordenatización que se reducirá módulo
tres. Por tanto, la identidad cuadrática y el código gráfico posterior
utilizan las mismas coordenadas, con sus incidencias y soportes.
